\section{Experimental Evaluation}

To validate the efficacy of ADRDS, a series of controlled experiments were conducted. This section outlines the experimental setup, the metrics analyzed, and the implications of the results for detection accuracy and system responsiveness.

\subsection{Experimental Setup}

All experiments were executed within a sandboxed environment using behavioral data generated by the ADRDS simulator. Ransomware-like event streams were synthesized alongside benign workloads mimicking standard user activity. The detection engine processed features in real time, with all activities logged via the interactive dashboard.

The experimental infrastructure comprised a dedicated workstation with an Intel Core i7 processor, 32GB RAM, and SSD storage. The sandboxed environment utilized Docker containers to ensure isolation and reproducibility. All experiments were repeated five times with different random seeds to assess result stability.

\subsection{Evaluation Metrics}

Performance was assessed using standard classification metrics—accuracy, precision, recall, and F$_1$-score—for the supervised classifiers. For the anomaly detector, reconstruction error distributions and threshold-dependent true positive and false positive rates were examined. Additionally, end-to-end latency was measured, defined as the interval from behavioral event generation to dashboard update, to assess the system's suitability for real-time interaction.

Accuracy measures the proportion of correct classifications among all samples. Precision quantifies the fraction of positive predictions that are actually positive, relevant for minimizing false alarms. Recall measures the fraction of actual positives correctly identified, critical for ensuring threats are not missed. F$_1$-score provides the harmonic mean of precision and recall, balancing both considerations.

For the anomaly detection component, we employed Receiver Operating Characteristic (ROC) curves and Area Under the Curve (AUC) metrics to evaluate performance across the full threshold range. Precision-Recall curves provided additional insight into performance under class imbalance.

\subsection{Experimental Parameters}

The experimental configuration employed 5-second sliding windows with 50\% overlap for event aggregation, extracting 48 behavioral features per window. The Random Forest classifier utilized 100 estimators with maximum depth 15 and minimum samples per leaf 5, while the LSTM autoencoder comprised two layers with 64 and 32 units respectively, a latent dimension of 16, and was trained for 50 epochs. Sensitivity analysis varied detection thresholds from 0.3 to 0.9 in 0.1 increments. The evaluation encompassed 20 distinct attack scenarios with encryption rates ranging from 100 to 10,000 files per minute across varied targeting strategies.

\subsection{Detection Performance}

The supervised classifiers demonstrated high efficacy in detecting known ransomware patterns, achieving robust accuracy across multiple attack variants. Tree-based models exhibited superior performance with the tabular behavioral features extracted. The Random Forest classifier achieved 94.3\% accuracy, 93.1\% precision, 95.2\% recall, and an F$_1$-score of 94.1\% on the held-out test set. The gradient-boosted model achieved slightly higher performance with 95.1\% accuracy.

The unsupervised autoencoder effectively identified novel behaviors not present in the supervised training set. Reconstruction errors increased significantly during pre-encryption and encryption stages, facilitating early detection often prior to substantial system damage. The AUC for the autoencoder reached 0.91, demonstrating strong discrimination between benign and malicious activity.

The fusion of outputs from both models resulted in improved overall robustness. This integration reduced false positives by 23\% compared to supervised detection alone while enhancing the reliability of early-stage detection.

\subsection{Latency and Responsiveness}

Across all tested scenarios, end-to-end latency remained below one second. Mean latency measured 340ms with a 95th percentile of 780ms. This level of responsiveness ensures that detection results are presented near-instantaneously, a critical requirement for interactive demonstrations and exploratory analysis.

Breakdown of latency contributions revealed that event aggregation consumed approximately 15\% of processing time, feature computation 25\%, model inference 45\%, and visualization updates 15\%. Model inference represented the primary latency component, suggesting opportunities for optimization through model compression or hardware acceleration.

\subsection{Ablation and Sensitivity Analysis}

Ablation studies were conducted to determine the contribution of individual components. Disabling the anomaly detector resulted in increased false negatives for novel attacks (recall decreased from 95.2\% to 87.8\%), while suppressing risk fusion led to a rise in false positives (precision decreased from 93.1\% to 84.6\%). Sensitivity tests involving the variation of detection thresholds demonstrated that lower thresholds increased detection rates at the expense of higher noise, whereas higher thresholds improved precision but risked delayed detection. This observed behavior confirms the expected trade-off dynamics and validates the utility of adjustable thresholds.

Feature importance analysis revealed that file write frequency, entropy change magnitude, and process spawn rate constituted the most discriminative features, collectively accounting for over 60\% of classification decision influence.

\subsection{Cross-Validation Methodology}

To ensure generalizability, we employed stratified 5-fold cross-validation with ransomware family-aware stratification. This approach ensured that samples from each ransomware family appeared in both training and test folds across different splits, preventing overfitting to specific family characteristics. Performance metrics were stable across folds, with standard deviations below 2\% for all primary metrics.

\subsection{Analysis Summary}

The experimental results confirm that ADRDS balances accuracy, interpretability, and processing speed. While the evaluation utilized simulated data, the use of models trained on established benchmarks and a structured taxonomy suggests generalizability to operational settings. Crucially, the system demonstrated the requisite speed and flexibility for interactive application.
